\documentclass{amsart}
% For dealing with references we use the comment environment
\usepackage{verbatim}
\newenvironment{reference}{\comment}{\endcomment}
%\newenvironment{reference}{}{}
\newenvironment{slogan}{\comment}{\endcomment}
\newenvironment{history}{\comment}{\endcomment}

% For commutative diagrams we use Xy-pic
\usepackage[all]{xy}

% We use 2cell for 2-commutative diagrams.
\xyoption{2cell}
\UseAllTwocells

% We use multicol for the list of chapters between chapters
\usepackage{multicol}

% This is generally recommended for better output
\usepackage{lmodern}
\usepackage[T1]{fontenc}

% For cross-file-references

% Package for hypertext links:
\usepackage{hyperref}

% For any local file, say "hello.tex" you want to link to please

% Theorem environments.
%
\theoremstyle{plain}
\newtheorem{theorem}[subsection]{Theorem}
\newtheorem{proposition}[subsection]{Proposition}
\newtheorem{lemma}[subsection]{Lemma}

\theoremstyle{definition}
\newtheorem{definition}[subsection]{Definition}
\newtheorem{example}[subsection]{Example}
\newtheorem{exercise}[subsection]{Exercise}
\newtheorem{situation}[subsection]{Situation}

\theoremstyle{remark}
\newtheorem{remark}[subsection]{Remark}
\newtheorem{remarks}[subsection]{Remarks}

\numberwithin{equation}{subsection}

% Macros
%
\def\lim{\mathop{\mathrm{lim}}\nolimits}
\def\colim{\mathop{\mathrm{colim}}\nolimits}
\def\Spec{\mathop{\mathrm{Spec}}}
\def\Hom{\mathop{\mathrm{Hom}}\nolimits}
\def\Ext{\mathop{\mathrm{Ext}}\nolimits}
\def\SheafHom{\mathop{\mathcal{H}\!\mathit{om}}\nolimits}
\def\SheafExt{\mathop{\mathcal{E}\!\mathit{xt}}\nolimits}
\def\Sch{\mathit{Sch}}
\def\Mor{\mathop{\mathrm{Mor}}\nolimits}
\def\Ob{\mathop{\mathrm{Ob}}\nolimits}
\def\Sh{\mathop{\mathit{Sh}}\nolimits}
\def\NL{\mathop{N\!L}\nolimits}
\def\CH{\mathop{\mathrm{CH}}\nolimits}
\def\proetale{{pro\text{-}\acute{e}tale}}
\def\etale{{\acute{e}tale}}
\def\QCoh{\mathit{QCoh}}
\def\Ker{\mathop{\mathrm{Ker}}}
\def\Im{\mathop{\mathrm{Im}}}
\def\Coker{\mathop{\mathrm{Coker}}}
\def\Coim{\mathop{\mathrm{Coim}}}

% Boxtimes
%
\DeclareMathSymbol{\boxtimes}{\mathbin}{AMSa}{"02}

%
% Macros for moduli stacks/spaces
%
\def\QCohstack{\mathcal{QC}\!\mathit{oh}}
\def\Cohstack{\mathcal{C}\!\mathit{oh}}
\def\Spacesstack{\mathcal{S}\!\mathit{paces}}
\def\Quotfunctor{\mathrm{Quot}}
\def\Hilbfunctor{\mathrm{Hilb}}
\def\Curvesstack{\mathcal{C}\!\mathit{urves}}
\def\Polarizedstack{\mathcal{P}\!\mathit{olarized}}
\def\Complexesstack{\mathcal{C}\!\mathit{omplexes}}
% \Pic is the operator that assigns to X its picard group, usage \Pic(X)
% \Picardstack_{X/B} denotes the Picard stack of X over B
% \Picardfunctor_{X/B} denotes the Picard functor of X over B
\def\Pic{\mathop{\mathrm{Pic}}\nolimits}
\def\Picardstack{\mathcal{P}\!\mathit{ic}}
\def\Picardfunctor{\mathrm{Pic}}
\def\Deformationcategory{\mathcal{D}\!\mathit{ef}}

\setlength{\emergencystretch}{3em}
\begin{document}
\title[Stacks Project, Tag 02TT]{450 --- Chow pullback is surjective for Zariski-locally trivial affine-space bundles}
\maketitle
\noindent Copyright (C) 2005--2025 Johan de Jong. Stacks Project authors. Source: \href{https://stacks.math.columbia.edu/tag/02TT}{Tag 02TT}.

\noindent Editorial difficulty note (not author text). Difficulty H2: The proof must reduce an arbitrary cycle to pullbacks while controlling locally finite image supports. Dimension induction and localization first reduce to a product affine line, where a generic irreducible polynomial makes every nonglobal component rationally trivial after shrinking and moves the remaining cycle to a lower-dimensional boundary..

\noindent Editorial provenance note (not author text). History: excerpt from revision \texttt{a0444\allowbreak 6e57e\allowbreak c1fbc\allowbreak 252a8\allowbreak 71afc\allowbreak ec775\allowbreak 2fb28\allowbreak 07b14}, chow.tex, lines 5650--5734. Reference presentation was adapted; original-author context and core support blocks were retained or added. The mathematical wording is unchanged. Licensed under GNU FDL 1.2 or later; no invariant sections or cover texts. See ../licenses/COPYING. Context blocks reproduce author definitions and supporting results; they are not additional counted problems.

\begin{situation}
\label{situation-setup}
Here $S$ is a locally Noetherian, and universally catenary scheme.
Moreover, we assume $S$ is endowed with a dimension function
$\delta : S \longrightarrow \mathbf{Z}$.
\end{situation}

\begin{lemma}
\label{lemma-pullback-affine-fibres-surjective}
Let $(S, \delta)$ be as in Situation \ref{situation-setup}.
Let $X$, $Y$ be locally of finite type over $S$.
Let $f : X \to Y$ be a flat morphism of relative dimension $r$.
Assume that for every $y \in Y$, there exists an open neighbourhood
$U \subset Y$ such that $f|_{f^{-1}(U)} : f^{-1}(U) \to U$
is identified with the morphism $U \times \mathbf{A}^r \to U$.
Then $f^* : \CH_k(Y) \to \CH_{k + r}(X)$ is surjective for all
$k \in \mathbf{Z}$.
\end{lemma}

\begin{proof}
Let $\alpha \in \CH_{k + r}(X)$.
Write $\alpha = \sum m_j[W_j]$ with $m_j \not = 0$ and
$W_j$ pairwise distinct integral closed subschemes of
$\delta$-dimension $k + r$. Then the family $\{W_j\}$
is locally finite in $X$. For any quasi-compact open
$V \subset Y$ we see that $f^{-1}(V) \cap W_j$
is nonempty only for finitely many $j$. Hence the
collection $Z_j = \overline{f(W_j)}$ of closures
of images is a locally finite collection of integral
closed subschemes of $Y$.

\medskip\noindent
Consider the fibre product diagrams
$$
\xymatrix{
f^{-1}(Z_j) \ar[r] \ar[d]_{f_j} & X \ar[d]^f \\
Z_j \ar[r] & Y
}
$$
Suppose that $[W_j] \in Z_{k + r}(f^{-1}(Z_j))$
is rationally equivalent to $f_j^*\beta_j$ for some
$k$-cycle $\beta_j \in \CH_k(Z_j)$. Then
$\beta = \sum m_j \beta_j$ will be a $k$-cycle on $Y$
and $f^*\beta = \sum m_j f_j^*\beta_j$ will be rationally
equivalent to $\alpha$ (see
Remark \href{https://stacks.math.columbia.edu/tag/02RZ}{remark infinite sums rational equivalences (Tag 02RZ)}).
This reduces us to the case $Y$ integral, and
$\alpha = [W]$ for some integral closed subscheme
of $X$ dominating $Y$. In particular we may
assume that $d = \dim_\delta(Y) < \infty$.

\medskip\noindent
Hence we can use induction on $d = \dim_\delta(Y)$.
If $d < k$, then $\CH_{k + r}(X) = 0$ and the lemma holds.
By assumption there exists a dense open $V \subset Y$ such
that $f^{-1}(V) \cong V \times \mathbf{A}^r$ as schemes over $V$.
Suppose that we can show that $\alpha|_{f^{-1}(V)} = f^*\beta$
for some $\beta \in Z_k(V)$. By Lemma \href{https://stacks.math.columbia.edu/tag/02RC}{lemma exact sequence open (Tag 02RC)}
we see that
$\beta = \beta'|_V$ for some $\beta' \in Z_k(Y)$.
By the exact sequence
$\CH_k(f^{-1}(Y \setminus V)) \to \CH_k(X) \to \CH_k(f^{-1}(V))$
of Lemma \href{https://stacks.math.columbia.edu/tag/02RX}{lemma restrict to open (Tag 02RX)}
we see that $\alpha - f^*\beta'$ comes from
a cycle $\alpha' \in \CH_{k + r}(f^{-1}(Y \setminus V))$.
Since $\dim_\delta(Y \setminus V) < d$ we win by
induction on $d$.

\medskip\noindent
Thus we may assume that $X = Y \times \mathbf{A}^r$.
In this case we can factor $f$ as
$$
X = Y \times \mathbf{A}^r \to
Y \times \mathbf{A}^{r - 1} \to \ldots \to
Y \times \mathbf{A}^1 \to Y.
$$
Hence it suffices to do the case $r = 1$. By the argument in the
second paragraph of the proof we are reduced to the case
$\alpha = [W]$, $Y$ integral, and $W \to Y$ dominant.
Again we can do induction on $d = \dim_\delta(Y)$.
If $W = Y \times \mathbf{A}^1$, then $[W] = f^*[Y]$.
Lastly, $W \subset Y \times \mathbf{A}^1$ is a proper inclusion,
then $W \to Y$ induces a finite field extension $R(W)/R(Y)$.
Let $P(T) \in R(Y)[T]$ be the monic irreducible polynomial such
that the generic fibre of $W \to Y$ is cut out by $P$ in
$\mathbf{A}^1_{R(Y)}$. Let $V \subset Y$ be a nonempty open such
that $P \in \Gamma(V, \mathcal{O}_Y)[T]$, and such that
$W \cap f^{-1}(V)$ is still cut out by $P$. Then we see that
$\alpha|_{f^{-1}(V)} \sim_{rat} 0$ and hence $\alpha \sim_{rat} \alpha'$
for some cycle $\alpha'$ on $(Y \setminus V) \times \mathbf{A}^1$.
By induction on the dimension we win.
\end{proof}
\end{document}
